\subsection{THETA3-C23: the actual A6 cycle has an integral model, with a precise KR comparison test}\label{C23}

\paragraph{Statement.}
Let $R$ be a discrete valuation ring with fraction field $L$
of characteristic zero and residue characteristic $p$.
Suppose the A6 data are defined over $L$, the smooth projective
$n$-fold $X$ has a smooth proper model $\mathcal X/R$, and
all the auxiliary abelian varieties $B_i,B_i'$ have abelian
scheme models $\mathcal B_i,\mathcal B_i'$ over $R$.
Write $\mathcal A'=\prod_{i=1}^{2n}\mathcal B_i'$ and
$g=\dim(\mathcal A'/R)$.
After the finite field extensions needed to define the finitely
many maps and divisor classes, the maps $f_i:X\to B_i$ and
all divisor kernels in A6 extend over $R$.

More precisely, write each selected pure cohomological divisor class as
\[
 D_i=\sum_j c_{ij}c_1(L_{ij}),\qquad c_{ij}\in M,
\]
for a number field $M$ and actual line bundles $L_{ij}$ on
$B_i\times B_i'$.
There are line bundles $\mathcal L_{ij}$ extending them and
maps $\mathfrak f_i:\mathcal X\to\mathcal B_i$ extending
$f_i$.  If
\[
 \pi:\mathcal X\times_R\mathcal A'\longrightarrow\mathcal A',
 \qquad
 g_i:\mathcal X\times_R\mathcal A'\longrightarrow
                   \mathcal B_i\times_R\mathcal B_i'
\]
are the projection and the map using $\mathfrak f_i$ and
the $i$th copied factor, put
\begin{equation}\label{C23-integral-cycle}
 \mathcal C=
  \pi_*\left(
    \prod_{i=1}^{2n}
        \left(\sum_jc_{ij}c_1(g_i^*\mathcal L_{ij})\right)
        \cap[\mathcal X\times_R\mathcal A']\right)
       \in\CH^n(\mathcal A')\otimes_\Z M.
\end{equation}
Its generic fibre is exactly the A6 output cycle for the chosen
divisor representatives.  Its refined special fibre is the same
explicit formula with every object reduced modulo $p$.
In particular, the A6 nonzero cohomology class remains nonzero
in prime-to-$p$ \'etale cohomology after specialization.

Here ``integral model'' refers to extension over $R$.
It does not remove the algebraic coefficients $c_{ij}$ or their
denominators.  One may clear them in a fixed coefficient field;
any resulting denominator at $p$ must be retained in a numerical
valuation comparison.

\paragraph{Why the maps extend at the specified good prime.}
An abelian scheme over a discrete valuation ring is the N\'eron
model of its generic fibre.  Its defining mapping property
therefore extends a generic-fibre map from the smooth
$R$-scheme $\mathcal X$ uniquely to $\mathcal B_i$.
It also extends every homomorphism between the auxiliary
abelian varieties uniquely.  These are applications at the
given prime, not spreading assertions obtained by deleting it.
An exact source is Conrad, \emph{Minimal models for elliptic
curves}, author version 21 November 2015, Definition~5.1 and
Example~5.3, p.~13; the latter records the same N\'eron property
for any abelian scheme.

The hypothesis that $\mathcal X$ is smooth remains substantive.
For the corresponding hyperspecial unitary moduli problem,
Li--Zhang, \S11.2, pp.~75--76, gives a smooth model under its
explicit hypotheses (H1)--(H6) and conditions (M1)--(M9).
One must match the selected component, prime and level to those
hypotheses; an arbitrary principal level at $p$ is not
hyperspecial.  Properness of the chosen model is also retained
in the statement above.

\paragraph{Why the divisor kernels extend, and their normalized form.}
For any smooth proper geometrically connected $R$-scheme
$Y$, the restriction map
\begin{equation}\label{C23-Pic}
                         \operatorname{Pic}(Y)
                           \longrightarrow\operatorname{Pic}(Y_L)
\end{equation}
is an isomorphism.  Here is the elementary divisor proof in
the situation used here.  The scheme $Y$ is regular and
integral, and its special fibre is smooth and geometrically
connected, hence geometrically integral.  A line bundle on
$Y_L$ has a nonzero rational section.  Close its finitely many
prime divisors in $Y$; their closures have codimension one.
Regular local rings are factorial, so this closed Weil divisor
is Cartier and defines an extension.  If an extended line
bundle is generically trivial, a generic trivialization gives
a Cartier divisor supported on the special fibre.  As that
fibre is integral, this divisor is an integer multiple of the
whole fibre, which is the principal divisor of a uniformizer.
Thus the line bundle was trivial.  This proves both assertions
in \eqref{C23-Pic}.  The factoriality input is Stacks Project,
Tag~0AG0; the rest of the argument has been given explicitly.
Apply it to $Y=\mathcal B_i\times_R\mathcal B_i'$.

The pure $H^1\otimes H^1$ kernel can, moreover, be represented
by normalized Poincar\'e correspondences.
For an individual line bundle $L$ on $B\times B'$, replace it by
\[
 L^{\rm norm}=L\otimes
       \operatorname{pr}_1^*(L|_{B\times0})^{-1}
       \otimes\operatorname{pr}_2^*(L|_{0\times B'})^{-1},
\]
with the constant fibre at $(0,0)$ rigidified.
Its first Chern class is the mixed K\"unneth component of
$c_1(L)$.  Since each $D_i$ is already purely mixed, this
operation leaves its cohomological class unchanged.
We may choose these normalized representatives before forming
\eqref{C23-integral-cycle}; its exact generic-fibre identity
then refers to that chosen cycle representative.  No equality
of unnormalized Chow classes is inferred from equality of
their cohomology classes.
The fibres of $L^{\rm norm}$ on $\{b\}\times B'$ are
algebraically trivial and the fibre at $b=0$ is trivial.
The universal property of the dual abelian variety yields
a map $h:B\to(B')^\vee$ with $h(0)=0$, hence a homomorphism,
and $L^{\rm norm}$ is the corresponding pullback of the
Poincar\'e bundle.  The exact universal property is Milne,
\emph{Abelian Varieties}, version~2.00 (2008), \S8,
p.~37, Remark~8.7(c); the zero-preserving homomorphism
property is \S16, Corollary~16.2(a), p.~67.

The homomorphism extends by the N\'eron property, and the
Poincar\'e kernel extends by \eqref{C23-Pic}; their compatibility
follows by uniqueness of the line-bundle extension.
Thus this construction does not require lifting an arbitrary
special-fibre divisor or inventing an integral Hodge projector.
It starts from the exact finite list of A6 correspondences.

\paragraph{Verification of the cycle and of specialization.}
The relative dimension of $\mathcal X\times_R\mathcal A'$
is $n+g$.  The product in \eqref{C23-integral-cycle} has
$2n$ first Chern operations, leaving relative dimension
$g-n$; proper pushforward therefore has codimension $n$
on $\mathcal A'$.
On the generic fibre, apply the projection formula to
$f\times1:X\times A'\to A\times A'$.
It identifies \eqref{C23-integral-cycle} with
\[
 p'_*(p^*(f_*[X])\cdot\prod_iq_i^*D_i),
\]
which is the actual A6 construction.  In the selected
cohomological component it is $(-1)^n k\,t_\beta$, with
the same nonzero $k$ and the same Koszul sign as A6.

The special-fibre immersion is the Cartier divisor defined
by a uniformizer.  Its refined Gysin map takes the fundamental
class of the smooth product to the special-fibre fundamental
class, commutes with first Chern operations, and commutes
with proper pushforward.  Consequently it takes
\eqref{C23-integral-cycle} to exactly the stated fibre formula.
Explicit algebraic references are Stacks Project,
Section~42.29, Tag~02T7, especially Lemma~42.29.8,
and Lemma~42.30.5, Tag~0B73, together with
Section~42.26, Tag~0AYA.  The original prime-to-$p$
cycle-class and specialization locators in C17 apply to this
formula.  Smooth proper base change is an isomorphism, so
the nonzero generic class cannot become zero.
CM endomorphism actions also extend, hence the specified
eigencomponent is preserved.  This does not say that the
subalgebra generated by special-fibre divisors has the same
dimension as in characteristic zero.

No proper intersection of chosen effective divisors is needed
to define \eqref{C23-integral-cycle}: first Chern operations
are defined for line bundles.  When a local intersection is
computed using equations, its appropriate derived Koszul
complex must still be used unless regularity of that sequence
has been established.

\paragraph{The exact common category: balanced Poincar\'e biextensions.}
There is a precise moduli interpretation of the normalized
kernels which sharpens the comparison problem.  Let $A_0,A$
be abelian schemes over a base $S$, with actions $\iota_0,\iota$
of the same commutative order $\mathcal O$.  Write $P_A$
for the Poincar\'e biextension on $A\times_S A^\vee$.
For a homomorphism $x:A_0\to A$, set
\[
                   L_x=(x\times1)^*P_A.
\]
This is a line bundle together with its rigidifications on
both zero axes and its compatible addition laws in both
variables.  Then $x\mapsto L_x$ is a natural equivalence
between homomorphisms $A_0\to A$ and isomorphism classes of
such biextensions on $A_0\times_S A^\vee$.
The exact primary source is SGA~7~I, Expos\'e~VIII,
\S\S3.1--3.2, printed pp.~254--255, in particular
\[
 \operatorname{Biext}^1(P,Q;\mathbb G_m)
            \simeq\Hom(P,Q^\vee)
                         \tag{VIII.3.2.2}
\]
and the symmetric formula~(VIII.3.2.4).
Apply this with $P=A_0,Q=A^\vee$ and use abelian biduality.
Section~3.1 identifies the same objects with divisorial
correspondences.  The construction by pullback of the
Poincar\'e bundle is compatible with every base change.
An automorphism of a biextension induces a bilinear map
$A_0\times A^\vee\to\mathbb G_m$, hence is trivial because
$\Hom(A^\vee,\mathbb G_m)=0$; the latter is exactly
formula~(VIII.3.2.1).  Thus the relevant isomorphisms, when
they exist, are unique.

Under this equivalence, $\mathcal O$-linearity of $x$ is
equivalent to the balanced identities
\begin{equation}\label{C23-balanced}
 (\iota_0(a)\times1)^*L_x
       \simeq(1\times\iota(a)^\vee)^*L_x
                         \qquad(a\in\mathcal O)
\end{equation}
as biextensions.  Indeed the left member corresponds to
$x\iota_0(a)$ and the right member corresponds to
$\iota(a)x$, by functoriality of the Poincar\'e bundle.
Here $\iota(a)^\vee$ is the actual dual endomorphism;
there is no additional conjugation in
\eqref{C23-balanced}.  Rosati conjugation enters when the
polarizations are used to form a Hermitian pairing.

Consequently, for a nilpotent immersion $S_0\subset S$
and a fixed $\mathcal O$-linear $x_0$ over $S_0$, the
following lifting problems are equivalent: lift $x_0$ as an
$\mathcal O$-linear homomorphism between the prescribed
abelian schemes, or lift $L_{x_0}$ as an $\mathcal O$-balanced
biextension, with both rigidifications and both addition
laws.  To prove this assertion, associate $L_x$ to a lift
$x$; conversely apply the displayed equivalence over $S$.
Its compatibility with restriction to $S_0$ and its
injectivity identify the resulting homomorphism with $x_0$
after restriction.  Equation~\eqref{C23-balanced} gives
the required linearity in both directions.  The structures
in this equivalence are exact; the refinement below explains
when they can be recovered from an actual line-bundle lift.
An equality of cohomology classes is not such a lift.

With polarizations $\lambda_0,\lambda$, the Gram entries
of a tuple are recovered from these kernels by recovering
their homomorphisms and taking
\[
           T_{ij}=\lambda_0^{-1}x_j^\vee\lambda x_i.
\]
Thus the condition of prescribed Gram matrix, and any
specified framing identifications, can be imposed on
either side without loss of information.  If the
semi-global convention uses $\Hom\otimes\Z_{(p)}$, clear
denominators by a specified integer prime to $p$ and retain
the corresponding Gram rescaling.  On the $p$-divisible
groups such multiplication is invertible, so it does not
alter the local lifting condition.  Multiplication by $p$
is not covered by this argument; a special
quasi-homomorphism with a $p$ denominator still has to meet
the original genuine-homomorphism extension condition.

This identifies an exact category shared by the normalized
A6 kernels and the global special-homomorphism moduli
problem.  To identify their cycles one must still realize
the auxiliary pair as the specified tautological PEL pair,
with the required dimensions, signatures, polarizations,
level and framing, and match its particular kernels.
The equivalence does not provide those identifications.
For an arbitrary normalized divisor representative,
membership in the balanced category for a specified pair of
$\mathcal O$-actions remains a condition to check.

For the rank-one CM factors of A6 this condition can be
arranged explicitly.  Let $B$ have CM by a field $E$,
with $H^1(B,\Q)$ of rank one over $E$, and choose an
$E$-compatible polarization $\lambda:B\to B^\vee$;
its Rosati involution on $E$ is complex conjugation.
Give $A=B^\vee$ the conjugate-dual action
\[
             \iota_A(a)=\iota_B(\bar a)^\vee.
\]
For every $a$ in the acting order, the homomorphism
$h_a=\lambda\iota_B(a):B\to A$ is $\mathcal O$-linear.
Indeed, for $b\in\mathcal O$, Rosati compatibility gives
\[
 h_a\iota_B(b)=\lambda\iota_B(ab)
      =\iota_B(\bar b)^\vee\lambda\iota_B(a)
      =\iota_A(b)h_a.
\]
Let $L_{h_a}$ be its normalized Poincar\'e kernel on
$B\times A^\vee=B\times B$.
It equals $(\iota_B(a)\times1)^*L_\lambda$.
The mixed polarization class $c_1(L_\lambda)$ pairs each
$E$-eigenline indexed by $\tau$ with the one indexed by
$\bar\tau$, and every such component is nonzero by
nondegeneracy of the polarization and the rank-one
hypothesis.

To isolate the first eigenline $\sigma$, write its
idempotent in $E\otimes_\Q\overline\Q$ as
$e_\sigma=\sum_j c_j a_j$, with $a_j$ in the acting
order and $c_j$ algebraic.  Then
\[
                     \sum_jc_jc_1(L_{h_{a_j}})
\]
is exactly the nonzero $\sigma\otimes\bar\sigma$
component of the polarization kernel.  One nonzero
algebraic rescaling makes it the selected pure divisor
class of A6.  Thus one can choose the actual finite
list of divisor representatives from balanced kernels,
with all coefficients and denominators explicit.
After defining the polarization and the finitely many
endomorphisms, each extends over the given good base.
The target action just constructed must still be matched
to the tautological PEL action and its signature;
the construction does not identify them automatically.

\paragraph{On the fixed infinitesimal PEL pair, a genuine line-bundle lift suffices.}
There is a further reduction for the actual deformation
problem.  Let $S$ be Artinian local, let $S_0\subset S$
be a nilpotent closed subscheme, and fix abelian schemes
$A_0,A$ over $S$ with their order actions.  Suppose
$x_0:(A_0)_{S_0}\to A_{S_0}$ is $\mathcal O$-linear.
Then $x_0$ lifts if and only if the actual line bundle
$L_{x_0}$ lifts to $A_0\times_S A^\vee$.
No chosen addition laws on the proposed line-bundle
lift are needed as extra input.

Here is the full proof.  A homomorphism lift gives a
line-bundle lift by Poincar\'e pullback.  Conversely,
normalize any proposed line-bundle lift $L$ by setting
\[
 N=L\otimes\operatorname{pr}_1^*(L|_{A_0\times0})^{-1}
       \otimes\operatorname{pr}_2^*(L|_{0\times A^\vee})^{-1}
       \otimes\pi^*(L|_{(0,0)}),
\]
where $\pi:A_0\times_S A^\vee\to S$.
This formula supplies compatible rigidifications on
both zero axes; it still reduces to the normalized
$L_{x_0}$.  SGA~7~I, Expos\'e~VII, Example~2.9.5,
printed p.~175 (following Corollary~2.9.4, p.~174),
states precisely that, for abelian schemes over any
base, the category of such birigidified line bundles
is equivalent to the category of biextensions.
Its argument uses the theorem of the square for the
two maps to the relative Picard functors in~(2.9.5.1).
Thus $N$ has the required unique biextension structure.
The classification in Expos\'e~VIII now produces a
homomorphism $x:A_0\to A$ reducing to $x_0$.

It remains to prove that $x$ is $\mathcal O$-linear;
this follows from elementary nilpotent rigidity.
If a morphism $h:B\to C$ of abelian schemes over
Artinian local $S$ reduces to zero over $S_0$,
its topological image lies at the zero point of the
special fibre.  Hence it factors through an affine
open neighbourhood of that point in $C$.
Since $H^0(B,\mathcal O_B)=H^0(S,\mathcal O_S)$,
the resulting map to this affine open factors through
$S$.  If $h$ preserves zero, it is therefore the zero
map.  Apply this to
\[
                h=x\iota_0(a)-\iota(a)x
                              \qquad(a\in\mathcal O).
\]
Its reduction is zero because $x_0$ is linear, and it
preserves zero.  Thus $x$ is linear for every $a$.
The same argument proves uniqueness of a homomorphism
lift.  This proves the asserted equivalence.

Consequently the KR obstruction below is also the
obstruction to lifting this specified normalized
Poincar\'e line bundle on the prescribed abelian
pair.  This statement concerns an actual line bundle
and a fixed infinitesimal moduli problem; neither its
first Chern class nor an unspecified auxiliary pair
can replace those data.

\paragraph{The actual KR lifting equation.}
The integral construction above removes an extension problem;
it does not yet impose the special-homomorphism moduli condition.
That condition has an exact local test.
For definiteness take the unramified quadratic local situation
over $\Q_p$ in Kudla--Rapoport I.
Let $A\to A_0=A/I$ be a square-zero thickening in the
nilpotent-$p$ deformation category, with compatible divided
powers on $I$.  Fix an ambient unitary deformation over $A$
whose reduction over $A_0$ belongs to $\mathcal Z(x)$.
Write $\mathbb D(\overline E)$ and $\mathbb D(X)$ for the
covariant Dieudonn\'e crystals evaluated on the thickening,
with Hodge filtrations $F_{\overline E}$ and $F_X$.
The crystalline map induced by the special homomorphism is
$\alpha_x:\mathbb D(\overline E)\to\mathbb D(X)$.
Its obstruction to lifting is
\begin{equation}\label{C23-obstruction}
 F_{\overline E,0}\xrightarrow{\alpha_x}
       \mathbb D(X)_0\longrightarrow(\operatorname{Lie}X)_0.
\end{equation}
Here the subscript $0$ labels the indicated
$\mathcal O_{K_p}$ eigenspace; it is not reduction to $A_0$.
The map reduces to zero over $A_0$.
Both its source and target have rank one, so after choosing
generators it is one element of $I$, and the lift exists
exactly when that element vanishes.
This is the explicit proof of Kudla--Rapoport,
\emph{Special cycles on unitary Shimura varieties I},
Proposition~3.5 and equations~(3.2)--(3.3), pp.~15--16
of the author preprint.  The proposition also proves that a
nonempty $\mathcal Z(x)$ is a relative Cartier divisor.
Li--Zhang, \S\S2.1--2.4, pp.~10--12, fixes the corresponding
polarization, signature, framing and derived-cycle conventions.

Thus a candidate comparison must supply a PEL moduli map
and actual special quasi-homomorphisms whose obstruction
equations match the intended divisors.  A homomorphism between
the fixed auxiliary good abelian schemes already extends on
the entire base.  Its lifting obstruction in that constant
family is identically zero.  Such a fixed-target homomorphism,
by itself, cannot define a nonempty proper KR divisor in a
varying unitary deformation space.

\paragraph{A sufficient equation-level comparison, with derived correction.}
There is a useful reduction of the required verification.
If $D,D'$ are closed subschemes of an $R$-flat scheme and
both are $R$-flat, then equality of their generic-fibre
closed subschemes implies $D=D'$ whenever they are compared
in the same ambient scheme.  Indeed their ideals are saturated
with respect to the uniformizer:
\[
 I_D=I_D[1/\varpi_R]\cap\mathcal O,
       \qquad
 I_{D'}=I_{D'}[1/\varpi_R]\cap\mathcal O,
\]
because the two quotient structure sheaves have no
$\varpi_R$-torsion.  The displayed equality follows immediately.

In particular, suppose a map from a regular comparison space
$T$ to the KR deformation space has been specified, and its
pullbacks of the selected KR divisors are $R$-flat Cartier
divisors.  If chosen flat Cartier representatives of the A6
kernel terms agree with these pullbacks on the generic fibre,
they agree integrally.  Locally their equations differ by units.
The corresponding two-term complexes
$[\mathcal O_T\xrightarrow{s_i}\mathcal O_T]$ are then
isomorphic, as are their derived tensor products.
This is an equality of supported perfect complexes, so it
retains every higher Tor term.
Any subsequently specified proper correspondence and finite
Euler characteristic respects this equality.
If the equations form a regular sequence, the Koszul
complex has only its ordinary quotient as homology; otherwise
the alternating higher-homology terms cannot be discarded.

The flatness condition is exact and necessary for this
generic-to-integral reduction.  In $R[[t]]$, the divisors
$(t)$ and $(\varpi_Rt)$ have the same generic divisor and
the same associated trivial line bundle.
Intersecting with the horizontal section $t=\varpi_R$
gives respectively the quotients $R/(\varpi_R)$ and
$R/(\varpi_R^2)$, of lengths one and two.
Both intersections are regular-sequence intersections, so
the discrepancy is already present before higher Tor occurs.
The second divisor has a vertical component and is not
$R$-flat.  Hence a generic identity or an identity of first
Chern classes does not determine the required arithmetic length.

\paragraph{A full-rank local KR cycle cannot be the horizontal A6 cycle.}
Let the local unitary space have rank $N$ and signature
$(1,N-1)$ in Li--Zhang's convention.
For a full-rank lattice $L$ of special quasi-homomorphisms,
$\mathcal Z(L)$ has empty horizontal part.
Li--Zhang, Lemma~2.10.1, p.~15, proves this directly:
a characteristic-zero lift would give an isogeny
$\overline E^{\,N}\to X$, forcing signature $(0,N)$,
contrary to $(1,N-1)$.
The derived full-rank cycle is consequently supported on
the special fibre.

Pullback, tensor product and proper pushforward over $R$
preserve the property of being generically zero.
Thus no correspondence made from those operations can carry
this full-rank vertical derived cycle to the A6 cycle with
its nonzero generic cohomology class.
A literal cycle comparison must instead use a horizontal
special-cycle stage (necessarily of rank less than $N$),
with a separately identified arithmetic intersection supplying
the final vertical contribution.
The full-rank density derivative computes that derived
arithmetic intersection; it is not itself an integral
representative of the horizontal A6 class.
\hfill$\square$

\paragraph{What remains to be constructed.}
At a prime satisfying the good-model hypotheses, the actual
A6 map and finite divisor kernel now have the explicit
integral realization \eqref{C23-integral-cycle}; this part of
the comparison is available.  The unresolved KR step is a
specific PEL moduli map and a generic divisor identification
whose pullbacks are flat, or a direct identity of the
supported perfect complexes defined by
\eqref{C23-obstruction}.  Those data would make the
derived correction and the normalization testable.
They are not consequences of the automorphic packet label,
the equality of Hodge classes, or smooth proper base change.
